%%%%%%%%%%%%%%%%%%%%%%%%%%%%%%%%%%%%%%%%%%%%%%%%%%%%%%%%%%%%%%%%%%%%%%%%%%%%%%%%
\section{Conclusion}
\label{sec:conclusion}
%%%%%%%%%%%%%%%%%%%%%%%%%%%%%%%%%%%%%%%%%%%%%%%%%%%%%%%%%%%%%%%%%%%%%%%%%%%%%%%%

This systematic literature review provides a structured assessment of the
research landscape surrounding \gls{dectnr}, covering 51 primary studies
published between 2020 and July 2026. Following a \gls{prisma}-based review
methodology, the study examined publication trends, technical research areas,
evaluation methodologies, experimental platforms, performance metrics,
application domains, and the principal research gaps affecting the continued
development of the technology.

The review shows that DECT NR+ has progressed from standardization and
simulation-oriented feasibility studies toward an increasingly diverse body of
research spanning PHY and MAC operation, multi-hop networking, scheduling and
resource allocation, reliability, energy efficiency, positioning, security,
coexistence, and heterogeneous network integration. PHY, MAC, and routing
currently constitute the most extensively investigated areas, while recent
work increasingly employs commercial hardware, SDR implementations, testbeds,
and application-oriented prototypes. Experimental studies in industrial
propagation, low-latency communication, professional audio, power-grid
monitoring, and measurement-driven scheduling demonstrate that DECT NR+ has
moved beyond purely conceptual evaluation toward practical implementation.

At the same time, the available evidence remains uneven across technical
domains and evaluation methods. Simulation and analytical modeling continue to
represent a substantial part of the literature, whereas large-scale physical
mesh experiments, heterogeneous traffic, mobility, interference, failure
conditions, and long-duration operation remain comparatively underexplored.
Positioning, security, coexistence, and multi-RAT integration are also less
mature than the core PHY, MAC, and routing research. Moreover, differences in
experimental configurations and reported \glspl{kpi} limit direct comparison
between studies.

A central conclusion of this review is that DECT NR+ performance should be
considered as a cross-layer and application-dependent property. PHY
configuration, channel access, scheduling, routing, topology, radio duty cycle,
reliability mechanisms, and hardware behavior interact in determining the
latency, reliability, throughput, energy consumption, and scalability observed
by an application. Consequently, no single configuration can be expected to
provide optimal performance across the diverse requirements of industrial
\gls{iot} applications.

Future research should therefore increasingly move from isolated optimization
of individual mechanisms toward application-oriented and experimentally
validated system design. Particular priorities include cross-layer scheduling
and resource allocation, deterministic and reliability-oriented communication,
energy-aware multi-hop operation, large-scale hardware validation, secure and
interoperable heterogeneous networking, and reproducible experimental
methodologies. Explicitly mapping application requirements to measurable
network \glspl{kpi} and configurable DECT NR+ mechanisms will be essential for
this transition.

Overall, the reviewed evidence establishes DECT NR+ as a promising platform for
decentralized and infrastructure-light industrial wireless communication, with
native multi-hop operation, flexible resource management, and a growing
experimental ecosystem. Its next stage of development will depend less on
demonstrating individual technical capabilities and more on integrating those
capabilities into predictable, scalable, energy-efficient, and experimentally
validated solutions for real industrial environments.